\section{Introduction}

A logical conclusion can remain unchanged when its premises are reordered, a relation is expressed through
its converse, or a language uses a different grammatical realization. A model that changes its answer under
such transformations presents a different evaluation problem from a model that consistently reaches the
wrong conclusion, and both differ from a model that appears stable only because it repeats a biased answer.
Aggregate accuracy on a single linguistic realization obscures all three distinctions: equal scores can
conceal failures on different instances, and high agreement can arise without any competence at all.
Behavioral tests and consistency evaluations therefore complement conventional task-level scores
\citep{ribeiro2020checklist, elazar2021pararel}.

Turkish offers an unusually well-controlled setting for making these distinctions measurable. It expresses
grammatical relations through productive suffixation and morphophonological alternations, and permits
constituent-order choices whose discourse effects a literal translation does not capture; computational
descriptions of Turkish morphology formalize the interaction between morphotactics and phonological
realization \citep{oflazer1994turkish}. These properties supply a family of \emph{answer-preserving}
transformations (premise reordering, converse constructions, constituent-order alternatives, and case or
possessive variants) that can be applied to the same underlying problem. We treat Turkish as a source
of controlled variation, not as a language presumed intrinsically harder for logical reasoning.

Existing resources do not answer how equivalent Turkish formulations affect transitive ordering inference.
General logical reasoning benchmarks \citep{yu2020reclor, liu2020logiqa, han2024folio} vary surface form
and content together; multilingual reasoning evaluations \citep{shi2023mgsm} compare independently sampled
items across languages; Turkish morphological-competence and minimal-pair benchmarks
\citep{ismayilzada2025morphology, basar2025turblimp, basar2026morphologyaware} intentionally manipulate
acceptability rather than holding the answer fixed; and paraphrase-consistency and metamorphic evaluations
\citep{elazar2021pararel, segura2016metamorphic} motivate, but do not provide, a matched bilingual ordering
design with exact answers. Competence-Based Evaluation (CBE) treats correctness across presentations of a
parameterized problem family as a behavioral requirement \citep{jiang2026cbe}; we adopt this framing and
extend it to Turkish--English matching, morphosyntactic interventions, lexical measurements, and
query-distance controls.

Turkish-CBE renders each latent strict total order as matched Turkish and English prompts and as controlled
Turkish presentation transformations, with exact symbolic answers, explicit entity-count and query-distance
factors, and split-level separation of latent orders. On top of this design, a correctness-aware analysis
decomposes paired outcomes on jointly valid support into robust correctness ($R$), shared-wrong agreement
($D$), and presentation dependence ($U$), reported together with raw invariance, a chance-corrected
agreement coefficient, valid-pair coverage, and strict all-planned robustness
(Section~\ref{sec:statistics}). These quantities are complementary: invariance alone cannot separate
correct consistency from stable error, and accuracy alone cannot reveal presentation dependence.

The framework is exercised by completed exploratory experiments. A registered pilot evaluates
Qwen2.5-0.5B-Instruct and \texttt{gpt-4o-2024-08-06} on matched prompts (Section~\ref{sec:pilot}); a
supplemental quantized local sweep evaluates four Qwen2.5 scales (0.5B--14B) on 16,618 matched prompts each
across all nine relation families (Section~\ref{sec:p1-local}); and a two-reviewer Turkish linguistic
assessment of 360 sampled stimuli was completed after these runs (Section~\ref{sec:review}); it found
unanimous semantic-equivalence judgments within the reviewed sample, identified one construction requiring
grammaticality-based exclusion, and documented contested naturalness.

This article makes three principal contributions:
\begin{enumerate}[leftmargin=1.6em]
  \item \textbf{A controlled bilingual reasoning benchmark.} A deterministic, reproducible framework that
  holds the latent problem, gold answer, names, and query geometry fixed while varying language and Turkish
  presentation, supporting more controlled comparisons than independently sampled multilingual test sets.
  The design does not claim to eliminate every linguistic or tokenization confounder; it makes the remaining
  ones explicit and measurable.
  \item \textbf{Correctness-aware presentation robustness.} A denominator-aware paired decomposition that
  separates correctness from response stability. Its components are related to existing agreement and
  paired-outcome methods rather than constituting a new statistical theory; their joint presentation is what
  carries the evaluation content. The constant-response behavior of the smaller evaluated models provides a
  concrete empirical demonstration that high response invariance does not indicate reasoning competence.
  \item \textbf{An empirical characterization of linguistic presentation sensitivity.} Within the evaluated
  models, datasets, and settings: model-scale differences in discriminative accuracy, English--Turkish
  performance differences, strong sensitivity to the P3 constituent-order presentation, complexity-dependent
  degradation, naming-condition effects, and the distinction between correct consistency and
  constant-response invariance. These are findings about the evaluated systems, not generalizations to all
  large language models.
\end{enumerate}

Turkish-CBE addresses five research questions. RQ1 concerns the paired Turkish--English accuracy
difference; RQ2 sensitivity to answer-preserving Turkish presentations; RQ3 the separation of entity count
from query distance and query type; RQ4 associations among relation vocabulary, corpus frequency,
tokenization, and predicted-answer bias; RQ5 whether reasoning prompting or optional adaptation improves
correct consistency as well as accuracy. The completed exploratory evidence bears on RQ1--RQ3 (and supplies
RQ4 diagnostics), while confirmatory answers to all five---and any RQ5 adaptation effect---are reserved for
the preregistered frozen-test design (Section~\ref{sec:future}).
